\documentclass[12pt]{article}
\usepackage[a4paper,margin=1in]{geometry}
\usepackage{iftex}
\ifXeTeX
\usepackage{fontspec}
\defaultfontfeatures{Ligatures=TeX}
\setmainfont{Times New Roman}
\else
\usepackage[T1]{fontenc}
\usepackage[utf8]{inputenc}
\usepackage{newtxtext}
\usepackage{newtxmath}
\fi
\usepackage{enumitem}
\usepackage{hyperref}
\usepackage{parskip}
\usepackage{microtype}
\usepackage[normalem]{ulem}
\usepackage{xcolor}

\newcommand{\practiceid}[1]{\texttt{#1}}
\newlength{\readinggutter}
\setlength{\readinggutter}{2.8em}
\newlength{\readingfirstindent}
\setlength{\readingfirstindent}{1.5em}
\newlength{\questionblockgap}
\setlength{\questionblockgap}{0.45\baselineskip}
\newlength{\exampleblockgap}
\setlength{\exampleblockgap}{0.65\baselineskip}
\newlength{\passageparagraphgap}
\setlength{\passageparagraphgap}{0.45\baselineskip}
\newenvironment{exampleblock}{\par\vspace{0.35\baselineskip}\begingroup\raggedright\setlength{\parskip}{0pt}\setlength{\parindent}{0pt}}{\par\endgroup\vspace{\exampleblockgap}}
\newcommand{\readinglineindent}[2]{%
\par\noindent%
\hangindent=\dimexpr\readinggutter+0.9em\relax%
\hangafter=1%
\makebox[\readinggutter][r]{\scriptsize\color{gray}#1}\hspace{0.9em}\hspace*{\readingfirstindent}#2%
\par}
\newcommand{\readingline}[2]{%
\par\noindent%
\hangindent=\dimexpr\readinggutter+0.9em\relax%
\hangafter=1%
\makebox[\readinggutter][r]{\scriptsize\color{gray}#1}\hspace{0.9em}#2%
\par}

\setlength{\parindent}{0pt}
\setlength{\parskip}{0.5em}
\setlist[enumerate]{topsep=0.25em,itemsep=0.18em,parsep=0pt,partopsep=0pt}
\setlength{\emergencystretch}{2em}

\begin{document}
\section*{TOEFL Practice 06 - Tryout}
Source of truth: DOCX only.
\section*{Listening Comprehension}
\subsection*{Part A}
DIRECTIONS: In Part A you will hear short conversations between two speakers. At the end of each conversation, a third speaker will ask a question about what the first two speakers said. Each conversation and each question will be spoken only one time. Therefore, you must listen carefully to understand what each speaker says. After you hear a conversation and the question, read the four choices and select the one that is the best answer to the question the speaker asked. Then, on your answer. sheet, find the number of the question and blacken the space that corresponds to the letter for the answer you have chosen. Blacken the space completely so that the letter inside the space does not show.\par
\begin{exampleblock}
Listen to the following example.\par
On the recording, you hear:\par
(Man) Does the car need to be filled?\par
(Woman) Mary stopped at the gas station on her way home.\par
(Narrator) What does the woman mean?\par
In your test book, you will read:\par
(A) Mary bought some food.\par
(B) Mary had car trouble.\par
(C) Mary went shopping.\par
(D) Mary bought some gas.\par
From the conversation you learn that Mary stopped at the gas station on her way home. The best answer to the question "Does the car need to be filled?" is (D), "Mary bought some gas." Therefore, the correct answer is (D).\par
Now let us begin Part A with question number 1.\par
\end{exampleblock}
\noindent\textbf{1.}
\begin{enumerate}[label=(\Alph*),leftmargin=*,itemsep=0.2em]
\item She needs to check her room.
\item She should leave the key at the desk.
\item She can leave any time the front desk is open.
\item She doesn't need to return the key
\end{enumerate}
\par
\vspace{0.25\baselineskip}
\noindent\textbf{2.}
\begin{enumerate}[label=(\Alph*),leftmargin=*,itemsep=0.2em]
\item He went to the movies yesterday.
\item He decided to go to the movies.
\item He was too tired to see the movie.
\item He stayed for the movie yesterday.
\end{enumerate}
\par
\vspace{0.25\baselineskip}
\noindent\textbf{3.}
\begin{enumerate}[label=(\Alph*),leftmargin=*,itemsep=0.2em]
\item The travel agent liked the book.
\item The book is about a travel agent.
\item The flight is full.
\item Travellers may have difficulty reading.
\end{enumerate}
\par
\vspace{0.25\baselineskip}
\noindent\textbf{4.}
\begin{enumerate}[label=(\Alph*),leftmargin=*,itemsep=0.2em]
\item She thinks that papers must be typed.
\item She believes that handwritten papers are acceptable.
\item Omitted words can be written in.
\item Typing skills are necessary for all students.
\end{enumerate}
\par
\vspace{0.25\baselineskip}
\noindent\textbf{5.}
\begin{enumerate}[label=(\Alph*),leftmargin=*,itemsep=0.2em]
\item He heard the noise.
\item He has decided already.
\item The man didn't like the noise.
\item The counsel will be noisy.
\end{enumerate}
\par
\vspace{0.25\baselineskip}
\noindent\textbf{6.}
\begin{enumerate}[label=(\Alph*),leftmargin=*,itemsep=0.2em]
\item Do you have the chemistry assignment?
\item Will Frank come to the chemistry lab?
\item Does Frank teach chemistry?
\item Will Frank be willing to help us?
\end{enumerate}
\par
\vspace{0.25\baselineskip}
\noindent\textbf{7.}
\begin{enumerate}[label=(\Alph*),leftmargin=*,itemsep=0.2em]
\item Only two pieces of luggage can be checked in
\item The boxes and the suitcases should be checked.
\item The suitcases belong to another passenger.
\item Suitcases cannot be checked in.
\end{enumerate}
\par
\vspace{0.25\baselineskip}
\noindent\textbf{8.}
\begin{enumerate}[label=(\Alph*),leftmargin=*,itemsep=0.2em]
\item Do you know where my coat is?
\item Have you been wearing my red coat?
\item Where is your coat?
\item Have you read it yet?
\end{enumerate}
\par
\vspace{0.25\baselineskip}
\noindent\textbf{9.}
\begin{enumerate}[label=(\Alph*),leftmargin=*,itemsep=0.2em]
\item He lives in the dorm.
\item He has a new apartment.
\item His apartment is north of campus.
\item He is moving soon.
\end{enumerate}
\par
\vspace{0.25\baselineskip}
\noindent\textbf{10.}
\begin{enumerate}[label=(\Alph*),leftmargin=*,itemsep=0.2em]
\item The wait is 15 minutes if you have a reservation.
\item Those who have to wait should make reservations.
\item The waiting room is on the fifth floor.
\item Without a reservation, you may wait for almost an hour,
\end{enumerate}
\par
\vspace{0.25\baselineskip}
\noindent\textbf{11.}
\begin{enumerate}[label=(\Alph*),leftmargin=*,itemsep=0.2em]
\item She thinks the food is excellent.
\item She feels that the price is reasonable.
\item She doesn't want to order the salad.
\item She thinks they should order dessert.
\end{enumerate}
\par
\vspace{0.25\baselineskip}
\noindent\textbf{12.}
\begin{enumerate}[label=(\Alph*),leftmargin=*,itemsep=0.2em]
\item Repair Jennifer's car.
\item Start driving to Jennifer's house.
\item Take Jennifer to the library.
\item Take Jennifer to the bank.
\end{enumerate}
\par
\vspace{0.25\baselineskip}
\noindent\textbf{13.}
\begin{enumerate}[label=(\Alph*),leftmargin=*,itemsep=0.2em]
\item He is excited about the show.
\item He has been to the museum many times.
\item He has seen the crafts displayed at the show.
\item He doesn't want to attend the craft show.
\end{enumerate}
\par
\vspace{0.25\baselineskip}
\noindent\textbf{14.}
\begin{enumerate}[label=(\Alph*),leftmargin=*,itemsep=0.2em]
\item Avoid traveling through the city.
\item Travel on a different day.
\item Stay away from highways.
\item Travel at 75 miles an hour.
\end{enumerate}
\par
\vspace{0.25\baselineskip}
\noindent\textbf{15.}
\begin{enumerate}[label=(\Alph*),leftmargin=*,itemsep=0.2em]
\item In the kitchen.
\item In a bookstore.
\item At a sports counter.
\item Near the outside door.
\end{enumerate}
\par
\vspace{0.25\baselineskip}
\noindent\textbf{16.}
\begin{enumerate}[label=(\Alph*),leftmargin=*,itemsep=0.2em]
\item She is the course advisor.
\item The man needs to take the history course.
\item The man is not taking a discussion course.
\item She cannot advise the man.
\end{enumerate}
\par
\vspace{0.25\baselineskip}
\noindent\textbf{17.}
\begin{enumerate}[label=(\Alph*),leftmargin=*,itemsep=0.2em]
\item It's around the corner.
\item It's in the next block.
\item It's on woman's right.
\item It's at the end of the block.
\end{enumerate}
\par
\vspace{0.25\baselineskip}
\noindent\textbf{18.}
\begin{enumerate}[label=(\Alph*),leftmargin=*,itemsep=0.2em]
\item At a packaging company.
\item In Seattle.
\item At a post office.
\item In a weight loss clinic.
\end{enumerate}
\par
\vspace{0.25\baselineskip}
\noindent\textbf{19.}
\begin{enumerate}[label=(\Alph*),leftmargin=*,itemsep=0.2em]
\item His wife does not like the man's brother.
\item He'll be busy on Saturday.
\item He'll be away for the weekend.
\item His wife will be away.
\end{enumerate}
\par
\vspace{0.25\baselineskip}
\noindent\textbf{20.}
\begin{enumerate}[label=(\Alph*),leftmargin=*,itemsep=0.2em]
\item A sales clerk.
\item A shoe-maker.
\item A receptionist.
\item A tailor.
\end{enumerate}
\par
\vspace{0.25\baselineskip}
\noindent\textbf{21.}
\begin{enumerate}[label=(\Alph*),leftmargin=*,itemsep=0.2em]
\item Go to the pool.
\item Make a phone call.
\item Correct the mistake.
\item Write a letter.
\end{enumerate}
\par
\vspace{0.25\baselineskip}
\noindent\textbf{22.}
\begin{enumerate}[label=(\Alph*),leftmargin=*,itemsep=0.2em]
\item It's been a great day.
\item We should do it all today.
\item We should stop for today.
\item Let's call Kay
\end{enumerate}
\par
\vspace{0.25\baselineskip}
\noindent\textbf{23.}
\begin{enumerate}[label=(\Alph*),leftmargin=*,itemsep=0.2em]
\item The lab is open every day.
\item The weather was terrible on Friday.
\item The test is at the end of the week.
\item The towels will be brought on Friday.
\end{enumerate}
\par
\vspace{0.25\baselineskip}
\noindent\textbf{24.}
\begin{enumerate}[label=(\Alph*),leftmargin=*,itemsep=0.2em]
\item In a garden
\item Near a pool
\item At an office
\item At a home
\end{enumerate}
\par
\vspace{0.25\baselineskip}
\noindent\textbf{25.}
\begin{enumerate}[label=(\Alph*),leftmargin=*,itemsep=0.2em]
\item Taking a class
\item Designing a building
\item Watching TV
\item Looking at pictures
\end{enumerate}
\par
\vspace{0.25\baselineskip}
\noindent\textbf{26.}
\begin{enumerate}[label=(\Alph*),leftmargin=*,itemsep=0.2em]
\item He doesn't know if the museum is open.
\item He has never been in a museum.
\item He thinks there is a museum nearby.
\item He doesn't have the directions.
\end{enumerate}
\par
\vspace{0.25\baselineskip}
\noindent\textbf{27.}
\begin{enumerate}[label=(\Alph*),leftmargin=*,itemsep=0.2em]
\item She has to take the boxes downstairs.
\item Would she mind helping him with the boxes?
\item The boxe's downstairs are heavy.
\item Did the packers bring all the boxes?
\end{enumerate}
\par
\vspace{0.25\baselineskip}
\noindent\textbf{28.}
\begin{enumerate}[label=(\Alph*),leftmargin=*,itemsep=0.2em]
\item He has changed jobs.
\item He has two children.
\item He has two jobs.
\item He is looking for a job.
\end{enumerate}
\par
\vspace{0.25\baselineskip}
\noindent\textbf{29.}
\begin{enumerate}[label=(\Alph*),leftmargin=*,itemsep=0.2em]
\item Marsha is having a difficult time.
\item Marsha is not very practical.
\item Rob's jokes are difficult to understand.
\item Rob's jokes are not funny.
\end{enumerate}
\par
\vspace{0.25\baselineskip}
\noindent\textbf{30.}
\begin{enumerate}[label=(\Alph*),leftmargin=*,itemsep=0.2em]
\item The store is closing.
\item The store is near the bay.
\item The employees get a good salary.
\item The store offers part-time jobs.
\end{enumerate}
\par
\vspace{0.25\baselineskip}
\subsection*{Part B}
DIRECTIONS: In this part of the test, you will hear longer conversations. After each conversation, you will hear several questions. The conversations and questions will not be repeated. After you hear a question, read the four possible answers in. your test book and select. The best answer. Then, on your answer sheet, find the number of the question and fill in the space that corresponds to the letter of the answer you have chosen. Remember, you are not allowed to take notes or write in your test book.\par
\begin{exampleblock}
Listen to the following example:\par
You will hear:\par
You will read:\par
(A) He has changed jobs.\par
(B) He has two children.\par
(C) He has two jobs.\par
(D) He is looking for a job.\par
From the conversation you learn that Tom has taken an additional job. The best answer to the question "Why is Tom tired?" is (C), "He has two jobs." Therefore the correct answer is (C).\par
\end{exampleblock}
\noindent\textbf{31.}
\begin{enumerate}[label=(\Alph*),leftmargin=*,itemsep=0.2em]
\item Various species of shark.
\item The oceans where sharks live.
\item Specific features of the shark's body.
\item Special characteristic of the shark's hearing.
\end{enumerate}
\par
\vspace{0.25\baselineskip}
\noindent\textbf{32.}
\begin{enumerate}[label=(\Alph*),leftmargin=*,itemsep=0.2em]
\item One or two
\item Approximately 45
\item About 100
\item More than 350
\end{enumerate}
\par
\vspace{0.25\baselineskip}
\noindent\textbf{33.}
\begin{enumerate}[label=(\Alph*),leftmargin=*,itemsep=0.2em]
\item Their attacks
\item Their teeth
\item Their hunting habits
\item Their electrical fields
\end{enumerate}
\par
\vspace{0.25\baselineskip}
\noindent\textbf{34.} eyesight.
\begin{enumerate}[label=(\Alph*),leftmargin=*,itemsep=0.2em]
\item They move slowly and quietly.
\item They don't turn over on their backs.
\item Their tails and gills.
\item They have sharp senses of hearing and
\end{enumerate}
\par
\vspace{0.25\baselineskip}
\noindent\textbf{35.}
\begin{enumerate}[label=(\Alph*),leftmargin=*,itemsep=0.2em]
\item He is new in the city.
\item He is sick.
\item His family needs to be treated.
\item His hospital doesn't provide the service.
\end{enumerate}
\par
\vspace{0.25\baselineskip}
\noindent\textbf{36.}
\begin{enumerate}[label=(\Alph*),leftmargin=*,itemsep=0.2em]
\item The man should go to the hospital.
\item The man needs to call the referral service.
\item The man's wife can make the phone calls.
\item The man needs to refer her to a physician.
\end{enumerate}
\par
\vspace{0.25\baselineskip}
\noindent\textbf{37.}
\begin{enumerate}[label=(\Alph*),leftmargin=*,itemsep=0.2em]
\item Whether a physician accepts new patients.
\item Whether patients like the physician.
\item Which doctors are employed by the referral service.
\item How often insurance claims need to be filed.
\end{enumerate}
\par
\vspace{0.25\baselineskip}
\noindent\textbf{38.}
\begin{enumerate}[label=(\Alph*),leftmargin=*,itemsep=0.2em]
\item Hospitals have dentists who require care.
\item She is not sure whether they are provided.
\item Dentists usually belong to an association.
\item Dental insurance is not accepted at hospitals.
\end{enumerate}
\par
\vspace{0.25\baselineskip}
\subsection*{Part C}
DIRECTIONS: In Part C you will hear short lectures and extended conversations. At the end of each, you will be asked several questions. Each lecture or conversation and each question will be spoken only one time. For this reason, you must listen carefully to understand what each speaker says. After you hear a question, read the four choices and select the one that best answers the question the speaker asked. Then, on your answer sheet, find the number of the question and blacken the space that contains the letter for the answer you have chosen. Answer all questions according to what is stated or implied in the lecture or conversation.\par
\begin{exampleblock}
Listen to this sample talk.\par
You will hear:\par
Now listen, to the following example.\par
You will hear:\par
You will read:\par
(A) By cars and carriages\par
(B) By bicycles, trains, and carriages\par
(C) On foot and by boat\par
(D) On board ships and trains\par
The best answer to the question "According to the speaker, how did people travel before the invention of the automobile?" is (B), "By bicycles, trains, and carriages." Therefore, the correct answer is (B). You are not allowed to make notes during the test.\par
\end{exampleblock}
\noindent\textbf{39.}
\begin{enumerate}[label=(\Alph*),leftmargin=*,itemsep=0.2em]
\item To discuss the importance of expensive clothes.
\item To explain the purpose of interviews.
\item To prepare listeners for a job interview.
\item To emphasize the importance of being polite.
\end{enumerate}
\par
\vspace{0.25\baselineskip}
\noindent\textbf{40.}
\begin{enumerate}[label=(\Alph*),leftmargin=*,itemsep=0.2em]
\item It should be expensive.
\item It can bring good fortune.
\item It is required of all interviewers..
\item It should be clean.
\end{enumerate}
\par
\vspace{0.25\baselineskip}
\noindent\textbf{41.}
\begin{enumerate}[label=(\Alph*),leftmargin=*,itemsep=0.2em]
\item They don't care about their appearance.
\item They don't arrive on time?
\item They borrow other people's clothes.
\item They lose money if they are careless.
\end{enumerate}
\par
\vspace{0.25\baselineskip}
\noindent\textbf{42.}
\begin{enumerate}[label=(\Alph*),leftmargin=*,itemsep=0.2em]
\item To match that of the interviewer.
\item Increase it to save time.
\item Slow it down and think before responding.
\item To meet the interviewer's expectations.
\end{enumerate}
\par
\vspace{0.25\baselineskip}
\noindent\textbf{43.}
\begin{enumerate}[label=(\Alph*),leftmargin=*,itemsep=0.2em]
\item A party in Boston.
\item A raid by American colonists.
\item Tea popularity in America.
\item The price of merchandise in Boston.
\end{enumerate}
\par
\vspace{0.25\baselineskip}
\noindent\textbf{44.}
\begin{enumerate}[label=(\Alph*),leftmargin=*,itemsep=0.2em]
\item In 1342
\item In 1660
\item In 1767
\item In 1773
\end{enumerate}
\par
\vspace{0.25\baselineskip}
\noindent\textbf{45.}
\begin{enumerate}[label=(\Alph*),leftmargin=*,itemsep=0.2em]
\item One
\item Two
\item Three
\item Four
\end{enumerate}
\par
\vspace{0.25\baselineskip}
\noindent\textbf{46.}
\begin{enumerate}[label=(\Alph*),leftmargin=*,itemsep=0.2em]
\item They wanted to profit from its sales.
\item They, wanted to destroy it.
\item The tea was not fit for consumption.
\item The tea was too expensive to buy
\end{enumerate}
\par
\vspace{0.25\baselineskip}
\noindent\textbf{47.}
\begin{enumerate}[label=(\Alph*),leftmargin=*,itemsep=0.2em]
\item 3 percent
\item 5 percent
\item 36 percent
\item 100 percent
\end{enumerate}
\par
\vspace{0.25\baselineskip}
\noindent\textbf{48.}
\begin{enumerate}[label=(\Alph*),leftmargin=*,itemsep=0.2em]
\item Events as they continue to occur.
\item Voters' attitudes at a particular time.
\item People's political beliefs in general.
\item Politician's views on current problems.
\end{enumerate}
\par
\vspace{0.25\baselineskip}
\noindent\textbf{49.}
\begin{enumerate}[label=(\Alph*),leftmargin=*,itemsep=0.2em]
\item They distributed the information.
\item They used a scientific method.
\item They reminded voters to go to the polls.
\item They made phone calls to ask for opinions.
\end{enumerate}
\par
\vspace{0.25\baselineskip}
\noindent\textbf{50.}
\begin{enumerate}[label=(\Alph*),leftmargin=*,itemsep=0.2em]
\item Roosevelt lost to Landon by a landslide.
\item The sample didn't include various groups of population.
\item Voters changed-their minds.
\item The sample of the population was not large enough.
\end{enumerate}
\par
\vspace{0.25\baselineskip}
\section*{Structure and Written Expression}
\subsection*{Sentence Completion}
This section is designed to test your ability to recognize language structures that are appropriate in standard written English. The questions in this section belong to two types, each of which has special directions.\par
DIRECTIONS: Questions 1-15 are partial sentences. Below each sentence you will see four words or phrases, marked (A), (B), (C), and (D). Select the one word or phrase that best completes the sentence. Then, on your answer sheet, find the number of the question and fill in the space that contains the letter for the answer you have chosen. Fill in the space completely.\par
\begin{exampleblock}
\textbf{EXAMPLE I}\par
\noindent Drying flowers is the best way......them.
\begin{enumerate}[label=(\Alph*),leftmargin=*,itemsep=0.2em]
\item to preserve
\item by preserving
\item preserve
\item preserved
\end{enumerate}
\vspace{0.25\baselineskip}
The sentence should state, "Drying flowers is the best way to preserve them." Therefore, the correct answer is (A).\par
\end{exampleblock}
\begin{exampleblock}
\textbf{EXAMPLE II}\par
\noindent Many American universities......as small, private colleges.
\begin{enumerate}[label=(\Alph*),leftmargin=*,itemsep=0.2em]
\item begun
\item beginning
\item began
\item for the beginning
\end{enumerate}
\vspace{0.25\baselineskip}
The sentence should state, "Many American universities began as small, private colleges." Therefore, the correct answer is (C). After you read the directions, begin work on the questions.\par
\end{exampleblock}
\noindent\textbf{1.} The red deer is a playful animal....... the American elk.
\begin{enumerate}[label=(\Alph*),leftmargin=*,itemsep=0.2em]
\item direct relationship
\item relative to the direction
\item directly related to
\item relating it to the direct
\end{enumerate}
\par
\vspace{0.25\baselineskip}
\noindent\textbf{2.} Sally Ride participated in launching Canadian communication satellites and.... pharmaceuticals.
\begin{enumerate}[label=(\Alph*),leftmargin=*,itemsep=0.2em]
\item conduct experiments with
\item experimented with conducting
\item conducting experiments with
\item experimental conduct with
\end{enumerate}
\par
\vspace{0.25\baselineskip}
\noindent\textbf{3.} Various regions of North America.... of European colonists, who spoke mostly English, Spanish, or French.
\begin{enumerate}[label=(\Alph*),leftmargin=*,itemsep=0.2em]
\item have retained the languages
\item retaining the languages
\item the languages have retained
\item the retaining of languages
\end{enumerate}
\par
\vspace{0.25\baselineskip}
\noindent\textbf{4.} Division means..... into smaller equal groups.
\begin{enumerate}[label=(\Alph*),leftmargin=*,itemsep=0.2em]
\item break up a number
\item a number is a break up
\item breaking up a number
\item broken up is a number
\end{enumerate}
\par
\vspace{0.25\baselineskip}
\noindent\textbf{5.} Market research investigation is the formulation of a problem that a client.....
\begin{enumerate}[label=(\Alph*),leftmargin=*,itemsep=0.2em]
\item seeks solution
\item to solve seeking
\item solving to seek
\item seeks to solve
\end{enumerate}
\par
\vspace{0.25\baselineskip}
\noindent\textbf{6.} In the early 18th century, Philadelphia and Boston were..... New York, which continued to grow rapidly in the next century.
\begin{enumerate}[label=(\Alph*),leftmargin=*,itemsep=0.2em]
\item largest and more prosperous than
\item larger and more prosperous than
\item prosperity and larger than
\item than larger and more prosperous
\end{enumerate}
\par
\vspace{0.25\baselineskip}
\noindent\textbf{7.} Farmers grow oats to feed cattle......
\begin{enumerate}[label=(\Alph*),leftmargin=*,itemsep=0.2em]
\item and sell to processing plants
\item after they are being sold to processing plants
\item and selling them to processing plants
\item process and sold to them to plants
\end{enumerate}
\par
\vspace{0.25\baselineskip}
\noindent\textbf{8.} Pigment mixed with water does not dissolve but remains....
\begin{enumerate}[label=(\Alph*),leftmargin=*,itemsep=0.2em]
\item suspending and forming paint
\item suspend and form to paint
\item suspended while painting
\item suspended to form paint
\end{enumerate}
\par
\vspace{0.25\baselineskip}
\noindent\textbf{9.} Ralph De Palma, a pioneer American race driver, preferred matches against..... to open competitions.
\begin{enumerate}[label=(\Alph*),leftmargin=*,itemsep=0.2em]
\item other driver
\item another driver
\item drivers other
\item driver, another
\end{enumerate}
\par
\vspace{0.25\baselineskip}
\noindent\textbf{10.} When one party arranges for a position in court, all involved....
\begin{enumerate}[label=(\Alph*),leftmargin=*,itemsep=0.2em]
\item notified and must
\item must be notified
\item notifying them must be
\item must be notifying
\end{enumerate}
\par
\vspace{0.25\baselineskip}
\noindent\textbf{11.} Alcohol, generally considered a depressant, decreases essential.....
\begin{enumerate}[label=(\Alph*),leftmargin=*,itemsep=0.2em]
\item functioning brain with
\item functioning as a brain
\item brain functions
\item brain as a function of
\end{enumerate}
\par
\vspace{0.25\baselineskip}
\noindent\textbf{12.} Cooperatives, owned by the people who use their services, sell goods produced by their members or purchased......
\begin{enumerate}[label=(\Alph*),leftmargin=*,itemsep=0.2em]
\item directly from farmers
\item direct to the farmers
\item farmers directly
\item farmers with direct
\end{enumerate}
\par
\vspace{0.25\baselineskip}
\noindent\textbf{13.} Marc Chagall's dreamlike paintings and.... come from his recollections of childhood images and mementoes.
\begin{enumerate}[label=(\Alph*),leftmargin=*,itemsep=0.2em]
\item brilliantly colored
\item brilliant coloring
\item brilliant color
\item brilliantly and color
\end{enumerate}
\par
\vspace{0.25\baselineskip}
\noindent\textbf{14.} Rationalism as a philosophical trend was based on the idea.....
\begin{enumerate}[label=(\Alph*),leftmargin=*,itemsep=0.2em]
\item reason is superior to experience that
\item that is reason superior to that experience
\item that reason is superior to experience
\item is that reason superior to experience
\end{enumerate}
\par
\vspace{0.25\baselineskip}
\noindent\textbf{15.} Rudolf Nureyev became recognized and admired for his stage activity and.....
\begin{enumerate}[label=(\Alph*),leftmargin=*,itemsep=0.2em]
\item impressive dance technique
\item technique dance is impressive
\item technique in dance impressive
\item is impressive in dance technique
\end{enumerate}
\par
\vspace{0.25\baselineskip}
\subsection*{Error Identification}
DIRECTIONS: In questions 16-40 every sentence has four words or phrases that are underlined. The four underlined portions of each sentence are marked (A), (B), (C), and (D). Identify the one word or phrase that makes the sentence incorrect. Then, on your answer sheet, find the number of the question and fill in the space that contains the letter for the answer you have chosen. Fill in the space completely.\par
\begin{exampleblock}
\textbf{EXAMPLE I}\par
Christopher Columbus \uline{has sailed}(A) from Europe in 1492 and \uline{discovered a}(B) \uline{new land}(C) he \uline{thought to}(D) be India.\par
The sentence should state, "Christopher Columbus sailed from Europe in 1492 and discovered a new land he thought to be India." Therefore, you should choose answer (A).\par
\end{exampleblock}
\begin{exampleblock}
\textbf{EXAMPLE II}\par
\uline{As the roles}(A) of people in society change, \uline{so does}(B) the \uline{rules of}(C) conduct in certain \uline{situations}(D).\par
The sentence should state, "As the roles of people in society change, so do the rules of\par
conduct in certain situations." Therefore, you should choose answer (B).\par
\end{exampleblock}
\noindent\textbf{16.} Balm \uline{which grown in}(A) \uline{damp and shady}(B) woodlands, \uline{is a tall}(C), fragrant herb \uline{of the mint}(D) species.
\par
\vspace{\questionblockgap}
\noindent\textbf{17.} Banks \uline{store cash}(A) in fireproof \uline{safes insured}(B) against fire, \uline{robberies}(C), floods, theft, and \uline{natural disaster}(D).
\par
\vspace{\questionblockgap}
\noindent\textbf{18.} During the Middle Ages, \uline{riddles were}(A) composed \uline{of poets}(B) and became \uline{a form}(C) of entertainment \uline{and art}(D).
\par
\vspace{\questionblockgap}
\noindent\textbf{19.} Rainwater carries \uline{unused chemicals}(A) from fields into \uline{streams or lakes}(B), where \uline{various compounds}(C) promote the \uline{rate of grown}(D) of weeds.
\par
\vspace{\questionblockgap}
\noindent\textbf{20.} Nuclear \uline{reactors produce}(A) energy \uline{by split}(B) the atom in the \uline{target material}(C) into two \uline{nearly equal}(D) parts.
\par
\vspace{\questionblockgap}
\noindent\textbf{21.} Clifford Odets' play Waiting for Lefty ranks \uline{among}(A) \uline{the most important}(B) works \uline{whom deal with}(C) the struggle of \uline{the working class}(D).
\par
\vspace{\questionblockgap}
\noindent\textbf{22.} \uline{To enjoy}(A) an opera fully, the listener \uline{should be}(B) familiar with the summary of the plot, \uline{particularly}(C) if the opera \uline{is singing}(D) in a foreign language.
\par
\vspace{\questionblockgap}
\noindent\textbf{23.} In 1968, the U.S. \uline{government effort}(A) \uline{to controls}(B) wages and prices \uline{to halt}(C) inflation \uline{had little}(D) effect.
\par
\vspace{\questionblockgap}
\noindent\textbf{24.} Annie Oakley \uline{learned}(A) to shoot \uline{at the age}(B) of eight and helped \uline{to support}(C) her family \uline{for hunting}(D) for a hotel restaurant.
\par
\vspace{\questionblockgap}
\noindent\textbf{25.} English is the \uline{official language}(A) of New Zealand and \uline{when spoken}(B) throughout \uline{the country}(C) while Maoris \uline{speak their own}(D) language.
\par
\vspace{\questionblockgap}
\noindent\textbf{26.} Cumulative rhymes repeat \uline{the informations}(A) and \uline{the wording}(B) in earlier verses that \uline{continue for}(C) seven \uline{to twelve cycles}(D).
\par
\vspace{\questionblockgap}
\noindent\textbf{27.} Animals \uline{inhabiting the sea}(A) bottom \uline{can be carried}(B) the \uline{weight of the water}(C) because \uline{it buoys them}(D).
\par
\vspace{\questionblockgap}
\noindent\textbf{28.} \uline{Making}(A) a photographic print \uline{required paper}(B) coated \uline{with light-sensitive}(C) chemicals that \uline{react to}(D) special solutions.
\par
\vspace{\questionblockgap}
\noindent\textbf{29.} Cabbage and lettuce have \uline{such a short}(A) stems and \uline{broad leaves}(B) that they \uline{appear to have}(C) \uline{no stems at all}(D).
\par
\vspace{\questionblockgap}
\noindent\textbf{30.} Plowing reduces \uline{the hard}(A) of the \uline{upper}(B) 6 to 16 inches of the \uline{earth's crust}(C) and allows air movement into the gaps \uline{between soil particles}(D).
\par
\vspace{\questionblockgap}
\noindent\textbf{31.} In basic terms, all communication is a \uline{process of exchange information}(A), \uline{imparting thoughts}(B) and ideas, and \uline{attempting to make}(C) oneself \uline{understood}(D) by listeners or readers.
\par
\vspace{\questionblockgap}
\noindent\textbf{32.} Iron and steel are rigid \uline{in their solid}(A) state and \uline{need to be}(B) \uline{melting}(C) when they are to be \uline{reshaped into}(D) new forms.
\par
\vspace{\questionblockgap}
\noindent\textbf{33.} Porous filters of charcoal \uline{can remove}(A) dust particles \uline{from the air}(B) if it \uline{is directed through}(C) them in \uline{a steadily stream}(D).
\par
\vspace{\questionblockgap}
\noindent\textbf{34.} The sender \uline{places stamps}(A) on packages \uline{as proof}(B) that postage \uline{for mailing}(C) an envelope or a package \uline{has paid}(D).
\par
\vspace{\questionblockgap}
\noindent\textbf{35.} Potash \uline{was originally}(A) obtained \uline{by running}(B) water through wood ashes and \uline{boiling the solution}(C) in \uline{cast-iron pot}(D).
\par
\vspace{\questionblockgap}
\noindent\textbf{36.} The earth's \uline{gravitational force}(A) for a given mass \uline{accords matter}(B) weight and \uline{diminish}(C) as an \uline{object moves}(D) away from the center of the earth.
\par
\vspace{\questionblockgap}
\noindent\textbf{37.} Allan Pinkerton organized \uline{groups of armed}(A) citizens \uline{whose services}(B) were available to \uline{employers for}(C) a daily fee \uline{to breaking labor}(D) strikes.
\par
\vspace{\questionblockgap}
\noindent\textbf{38.} The Bureau of Labor Statistics \uline{found that}(A) \uline{few as}(B) 14 percent of \uline{executives}(C) write \uline{their own}(D) memos, letters, and speeches.
\par
\vspace{\questionblockgap}
\noindent\textbf{39.} Real estate law \uline{holds that}(A) \uline{not one}(B) may legally plant or remove plants \uline{of any kind}(C) without \uline{permission of}(D) the land owner.
\par
\vspace{\questionblockgap}
\noindent\textbf{40.} National income \uline{is defined}(A) as \uline{the total income}(B) \uline{earned, but}(C) not necessarily received by all \uline{persons lived}(D) in a country during a period of time.
\par
\vspace{\questionblockgap}
\section*{Reading Comprehension}
\begin{center}
\textbf{Time 55 minutes}
\end{center}
DIRECTIONS: In this section you will read several passages. Each passage is followed by a series of questions. For questions 1-50, you need to select the best answer, (A), (B), (C), or (D), to each question. Then, on your answer sheet, find the number of the question and fill in the space that contains the letter of the answer you have selected. Fill in the space completely. Answer all questions following a passage on the basis of what is stated or implied in the passage.\par
\begin{exampleblock}
\small
\sloppy
Read the following passage:\par
A tomahawk is a small ax used as a tool and a weapon by the North American Indian\par
tribes. An average tomahawk was not very long and did not weigh a great deal. Originally,\par
the head of the tomahawk was made of a shaped stone or an animal bone and was mounted on\par
Line a wooden handle. After the arrival of the European settlers, the Indians began to use toma-\par
hawks with iron heads. Indian males and females of all ages used tomahawks to chop and cut wood, pound stakes into the ground to put up wigwams, and do many other chores. Indian warriors relied on tomahawks as weapons and even threw them at their enemies. Some types of tomahawks were used in religious ceremonies. Contemporary American idioms reflect\par
this aspect of American heritage.\par
\end{exampleblock}
\begin{exampleblock}
\textbf{EXAMPLE I}\par
\noindent Early tomahawk heads were made of
\begin{enumerate}[label=(\Alph*),leftmargin=*,itemsep=0.2em]
\item stone or bone
\item wood or sticks
\item European iron
\item religious weapons
\end{enumerate}
\vspace{0.25\baselineskip}
According to the passage, early tomahawk heads were made of stone or bone. Therefore, the correct answer is (A).\par
\end{exampleblock}
\begin{exampleblock}
\textbf{EXAMPLE II}\par
\noindent How has the Indian use of tomahawks affected American daily life today?
\begin{enumerate}[label=(\Alph*),leftmargin=*,itemsep=0.2em]
\item Tomahawks are still used as weapons.
\item Tomahawks are used as tools for.certain jobs.
\item Contemporary language refers to tomahawks.
\item Indian tribes cherish tomahawks as heirlooms.
\end{enumerate}
\vspace{0.25\baselineskip}
The passage states, "Contemporary American idioms reflect this aspect of American heritage." The correct answer is (C). After you read the directions, begin work on the questions.\par
\end{exampleblock}
\subsection*{Questions 1-10}
\begingroup
\small
\sloppy
\setlength{\parskip}{0pt}
\setlength{\parindent}{0pt}
Free-standing sculpture that is molded or carved is a type familiar to almost everyone.\par
Although certain free-standing figures or groups of figures can have only a single side\par
intended for viewing, others are completed on all sides. As with all 'other forms of art, the\par
ultimate shape of a sculpture reflects the artist's vision of individuals or experiences repre-\par
sented by the work. Throughout history, people everywhere have discovered a need for\par
sculpture as a record of events and feelings.\par
Materials which can be sculpted do much to contribute to the artist's imagination.\par
Wood, stone, metal, and various types of plastic and synthetics are all used as sculpting media. When sculptures are made of stone, wood, ivory, or even ice, the sculptor carves or chips\par
the substance to reduce it to the necessary shape. Developing a sculptured image on all sides\par
represents a change from the older approach when artists left the back portion of the figure\par
unfinished and rough. In fact, sculpture in relief is completely attached to the flat background material and appears to be a part of it. Relief, which is completed only on one side intended for viewing, was the first type of sculpture created by man, when ancient sculptors\par
removed the background material in a side of a tree or a cave to make their drawing appear\par
more realistic.\par
While creating a statue, the artist 'depends on the appropriate lighting to develop the figure because the quality of the final product relies on the interplay between light and shade.\par
When the work is finished, the' sculpture must be displayed in the same light as it was origi-\par
nally created. If a light from a source is too weak or too strong, the effect that the sculptor intended may be lost. For example, in painting, the light and shade give the image shape and\par
solidity that cannot be altered by an external light in which it is displayed. When a sculpture is exhibited, the artist's work is brought to life by light, and its character can be altered\par
by the control of the light source. A fundamental difference betweena painting and a sculp-\par
ture is that when viewing a painting, the audience can only see the point of view that the\par
painter had intended. A free-standing sculpture can be seen from practically any angle. The\par
job of the sculptor is then to attain the quality and the volume of the image from any possible point of view.\par
In addition to carving a work, sculptures can be cast. In the process of casting, a sculpture\par
can be reproduced in a mold when a liquefied medium is poured intoa shape. After the material from which the sculpture is made hardens, the mold is removed, and the work is\par
cleaned of the excess and polished. Casting allows the artist to produce as many replicas as\par
needed. Most commercially sold sculptures are made in this way. Casting metals requires\par
special care and skill. Bronze is the preferred metal because of its versatility and malleabil-\par
ity. To make bronze sculpture, the space in a mold is filled with wax until it is melted by the\par
heated metal. This process, sometimes called lost-wax, was favored by Benvenuto Cellini\par
and was common among the artists in ancient China.\par
\endgroup
\medskip
\noindent\textbf{1.} What is the main topic of this passage?
\begin{enumerate}[label=(\Alph*),leftmargin=*,itemsep=0.2em]
\item Differences between painting and sculpture
\item Sculpting techniques and media
\item Types of commercially produced sculptures
\item Reasons for enjoying sculpture
\end{enumerate}
\par
\vspace{0.25\baselineskip}
\noindent\textbf{2.} The word "ultimate" in line 4 is closest in meaning to
\begin{enumerate}[label=(\Alph*),leftmargin=*,itemsep=0.2em]
\item ulterior
\item final
\item formal
\item formidable
\end{enumerate}
\par
\vspace{0.25\baselineskip}
\noindent\textbf{3.} According to the passage, the purpose of sculpture as a form of art is to
\begin{enumerate}[label=(\Alph*),leftmargin=*,itemsep=0.2em]
\item Display a group of figures
\item Reflect a human need for freedom
\item Express an artistic vision
\item Commemorate individuals and events
\end{enumerate}
\par
\vspace{0.25\baselineskip}
\noindent\textbf{4.} According to the passage, all of the following are true of sculpture EXCEPT that
\begin{enumerate}[label=(\Alph*),leftmargin=*,itemsep=0.2em]
\item it can be found in all parts of the world
\item it has undergone change since the early times
\item it can be created from many substances
\item it is no longer useful for people
\end{enumerate}
\par
\vspace{0.25\baselineskip}
\noindent\textbf{5.} The author of the passage implies that the most important factor in showing a sculpted work is
\begin{enumerate}[label=(\Alph*),leftmargin=*,itemsep=0.2em]
\item the strength of the light source
\item the development of the sculpted figure
\item the shape of the material for sculpting
\item the effect of light on the sculpted image
\end{enumerate}
\par
\vspace{0.25\baselineskip}
\noindent\textbf{6.} The word "audience" in line 19 is closest in meaning to
\begin{enumerate}[label=(\Alph*),leftmargin=*,itemsep=0.2em]
\item listeners
\item viewers
\item public
\item artists
\end{enumerate}
\par
\vspace{0.25\baselineskip}
\noindent\textbf{7.} What does the author mention as an important difference between a painting and a sculpture?
\begin{enumerate}[label=(\Alph*),leftmargin=*,itemsep=0.2em]
\item A painting does not need shading to be displayed.
\item A painting can be viewed from only one position.
\item A sculpture needs to have proper light.
\item A sculpture does not look good from all angles.
\end{enumerate}
\par
\vspace{0.25\baselineskip}
\noindent\textbf{8.} Which of the following is NOT mentioned as a sculpting medium?
\begin{enumerate}[label=(\Alph*),leftmargin=*,itemsep=0.2em]
\item Ice
\item Ivory
\item Stone
\item Wax
\end{enumerate}
\par
\vspace{0.25\baselineskip}
\noindent\textbf{9.} The word "replicas" in line 24 is closest in meaning to
\begin{enumerate}[label=(\Alph*),leftmargin=*,itemsep=0.2em]
\item replacements
\item molds
\item reproductions
\item monuments
\end{enumerate}
\par
\vspace{0.25\baselineskip}
\noindent\textbf{10.} According to the passage, what are the two basic methods for making sculptures?
\begin{enumerate}[label=(\Alph*),leftmargin=*,itemsep=0.2em]
\item Carving and casting
\item Free-standing and relief
\item Hardening and melting
\item Stone and metal
\end{enumerate}
\par
\vspace{0.25\baselineskip}
\subsection*{Questions 11-21}
\begingroup
\small
\sloppy
\setlength{\parskip}{0pt}
\setlength{\parindent}{0pt}
The Beatles became the most popular group in rock music history. This quartet of extraordinarily talented musicians generated a phenomenal number of pieces that won gold\par
records. They inspired a frenzy that transcended countries and economic strata. While all of\par
them sang, John Lennon and Paul McCartney wrote the majority of their songs. Originally,\par
Lennon and five others formed a group called the Quarrymen in 1956, with McCartney joining them later that year. George Harrison, John Lennon, and Paul McCartney, together\par
with Stuart Sutcliffe, who played the bass guitar, and Pete Best on the drums, performed together in several bands for a few years, until they finally settled on the Silver Beatles in\par
1960. American rock musicians, such as Chuck Berry and Elvis Presley, influenced Len-\par
non's and McCartney's music, whose first hits consisted of simple tunes and lyrics about\par
young love, "Love Me Do" and "Please, Please Me." The Beatles' U.S. tour propelled them\par
to stardom and led to two movies A Hard Day's Night and Help!, filmed in 1964 and 1965.\par
The so-called British invasion of the United States was in full swing when they took the top\par
five spots on the singles charts, followed by the release of their first film.\par
During the 1960s, their music matured and acquired a sense of melody. The lyrics of\par
their songs became deeper and gained in both imagination and meaning. Their popularity\par
continued to grow as the Beatles turned their attention to social problems and political issues in "Nowhere Man" and "Eleanor Rigby." Loneliness and nostalgia come through in\par
their ballads "Michelle" and "Yesterday," which fully displayed the group's professional de-\par
velopment and sophistication. Lennon's sardonic music with lyrics written in the first person, and McCartney's songs that created scenarios with off beat individuals, contributed to\par
the character of the music produced by the group. In addition to their music, the Beatles set\par
a social trend that popularized long hair, Indian music, and mod dress.\par
For a variety of reasons, the musicians began to drift apart, and their last concert took\par
place in San Francisco in 1966. The newspapers and tabloids publicized their quarrels and\par
lawsuits, and the much idolized group finally disbanded in 1970. However, their albums\par
had outsold those of any other band in history. Although all of the Beatles continued to perform solo or form new rock groups, alone, none could achieve the recognition and success\par
that they had been able to win together.\par
\endgroup
\medskip
\noindent\textbf{11.} What does the passage mainly discuss?
\begin{enumerate}[label=(\Alph*),leftmargin=*,itemsep=0.2em]
\item The history and music of the Beatles
\item The history and milestones of rock
\item The fashion and music popular in the 1960s
\item The creation and history of a music group
\end{enumerate}
\par
\vspace{0.25\baselineskip}
\noindent\textbf{12.} According to the passage, how many members were in the band, formed in 1956?
\begin{enumerate}[label=(\Alph*),leftmargin=*,itemsep=0.2em]
\item Four
\item Five
\item Six
\item Seven
\end{enumerate}
\par
\vspace{0.25\baselineskip}
\noindent\textbf{13.} According to the passage, which of the Beatles had the greatest musical talent?
\begin{enumerate}[label=(\Alph*),leftmargin=*,itemsep=0.2em]
\item John Lennon and Paul McCartney
\item George Harrison and John Lennon
\item Stuart Sutcliffe and Pete Best
\item John Lennon, Paul McCartney, and George Harrison
\end{enumerate}
\par
\vspace{0.25\baselineskip}
\noindent\textbf{14.} The author of the passage implies that the Beatles
\begin{enumerate}[label=(\Alph*),leftmargin=*,itemsep=0.2em]
\item competed with American musicians
\item wrote their music as a group
\item became popular relatively quickly
\item were active in social movements
\end{enumerate}
\par
\vspace{0.25\baselineskip}
\noindent\textbf{15.} According to the passage, the Beatles' fame grew as a result of
\begin{enumerate}[label=(\Alph*),leftmargin=*,itemsep=0.2em]
\item Chuck Berry's involvement
\item their American tour
\item two movies made in the U.S.
\item their first two hits
\end{enumerate}
\par
\vspace{0.25\baselineskip}
\noindent\textbf{16.} The author of the passage implies that over time, the music and lyrics by the Beatles
\begin{enumerate}[label=(\Alph*),leftmargin=*,itemsep=0.2em]
\item became more complex than at the beginning of their career
\item declined in quality and political significance
\item were dedicated to women named Eleanor and Michelle
\item made them the richest musicians in the world
\end{enumerate}
\par
\vspace{0.25\baselineskip}
\noindent\textbf{17.} The word "acquired" in line 12 is closest in meaning to
\begin{enumerate}[label=(\Alph*),leftmargin=*,itemsep=0.2em]
\item imparted
\item attached
\item imprinted
\item attained
\end{enumerate}
\par
\vspace{0.25\baselineskip}
\noindent\textbf{18.} According to the passage, when did the Beatles experience their greatest success?
\begin{enumerate}[label=(\Alph*),leftmargin=*,itemsep=0.2em]
\item In the late 1950s.
\item After their break-up in 1970.
\item During the early and mid-1960s.
\item Throughout their lifetimes.
\end{enumerate}
\par
\vspace{0.25\baselineskip}
\noindent\textbf{19.} The word "scenarios" in line 16 is closest in meaning to
\begin{enumerate}[label=(\Alph*),leftmargin=*,itemsep=0.2em]
\item sceneries
\item situations
\item life stories
\item love themes
\end{enumerate}
\par
\vspace{0.25\baselineskip}
\noindent\textbf{20.} According to the passage, how did Lennon and McCartney enhance the music of the group?
\begin{enumerate}[label=(\Alph*),leftmargin=*,itemsep=0.2em]
\item They struggled to reach stardom in the United States.
\item They composed lyrics to scornful songs and ballads.
\item Their music added distinctiveness to the Beatles' repertoire.
\item Their loneliness and sadness made their music popular.
\end{enumerate}
\par
\vspace{0.25\baselineskip}
\noindent\textbf{21.} In line 21, the word "disbanded" is closest in meaning to
\begin{enumerate}[label=(\Alph*),leftmargin=*,itemsep=0.2em]
\item separated
\item slipped
\item revelled
\item bonded
\end{enumerate}
\par
\vspace{0.25\baselineskip}
\subsection*{Questions 22-31}
\begingroup
\small
\sloppy
\setlength{\parskip}{0pt}
\setlength{\parindent}{0pt}
Like Europeans who amived in the Americas, the first American Indians were 'immigrants. Because\par
Indians were nomadic hunters and gatherers, they probably arrived in search of new hunting grounds from Asia when they crossed the ice-covered Bering Strait to Alaska. Anthropologists estimate that the entire Indian population north of Mexico was slightly greater\par
than 1,020,000 when the first settlers arrived from Europe. Although Native Americans belonged to one geographic race, their cultures and languages were only marginally similar,\par
and by and large, they had different ways of life. Nomadic migrations required Indians to\par
construct shelters that did not need to be transported, but could be easily erected from the\par
materials found in their new location.\par
Eastern Woodland Indian tribes lived in bark-covered wigwams that were shaped like\par
cones or domes. The frame for the hut was made of young trees firmly driven into the\par
ground, and then bent overhead to tie together with bark fibers or strings of animal hides.\par
Sheets and slabs of bark were attached to the frame to construct the roof and walls, leaving\par
an opening to serve as a door and to allow smoke to escape. The Iroquois in north eastern re-\par
gions built longhouses that were more spacious than wigwams because five to a dozen families lived under one roof. During the winter, they plastered clay to the poles of the frame to\par
protect the inhabitants from wind and rain.\par
Pueblo Indians who lived in the southwest portion of the United States in northern Arizona and New Mexico constructed elaborate housing with several stories and many rooms.\par
Each family unit had only one room, and their ancestors dug shelters in the walls of cliffs and\par
canyons. The ground story of a Pueblo dwelling had no doors or windows in order to prevent\par
enemies from entering. The next level was set back the width of one room, and the row of\par
rooms above it was set back once again, giving their houses the appearance of a terrace. Pueblos used ladders to climb to the upper levels and pulled them in when all family members re-\par
turned for the night.\par
Indians living in deserts used sandstone and clay as construction materials. Those who\par
lived in the valleys of rivers even made bricks of clay with wood chips to add strength and to\par
prevent the clay from cracking. To make roofs, Pueblos tied logs together to make rafters and\par
laid them across the two outside walls. On top of the rafters, layers of tree branches, sticks,\par
grass, and brush created a solid roof to preclude the water from leaking inside. Pueblo dwellings were dark because windows were often not large enough to allow much light.\par
\endgroup
\medskip
\noindent\textbf{22.} What does the passage mainly discuss?
\begin{enumerate}[label=(\Alph*),leftmargin=*,itemsep=0.2em]
\item Different Indian tribes
\item Types of households among Indians
\item Types of shelters built by Indians
\item Different Indian cultures
\end{enumerate}
\par
\vspace{0.25\baselineskip}
\noindent\textbf{23.} In line 3, the word "marginally" is closest in meaning to
\begin{enumerate}[label=(\Alph*),leftmargin=*,itemsep=0.2em]
\item markedly
\item minimally
\item temporarily
\item tentatively
\end{enumerate}
\par
\vspace{0.25\baselineskip}
\noindent\textbf{24.} In line 4, the phrase "by and large" is closest in meaning to
\begin{enumerate}[label=(\Alph*),leftmargin=*,itemsep=0.2em]
\item mostly
\item conversely
\item occasionally
\item notably
\end{enumerate}
\par
\vspace{0.25\baselineskip}
\noindent\textbf{25.} The author of the passage implies that Indians
\begin{enumerate}[label=(\Alph*),leftmargin=*,itemsep=0.2em]
\item carried their construction materials to new locations
\item lived in settlements, similar to the Europeans
\item liked the climate in the Southwest
\item constructed shelters every time a tribe moved
\end{enumerate}
\par
\vspace{0.25\baselineskip}
\noindent\textbf{26.} According to the passage, what shape did the shelters of Woodland Indians have?
\begin{enumerate}[label=(\Alph*),leftmargin=*,itemsep=0.2em]
\item random
\item round
\item rectangular
\item convex
\end{enumerate}
\par
\vspace{0.25\baselineskip}
\noindent\textbf{27.} The author of the passage implies that Eastern Indians
\begin{enumerate}[label=(\Alph*),leftmargin=*,itemsep=0.2em]
\item constructed huts without roofs
\item planted trees to harvest crops
\item made fires inside their huts
\item used sheets and blankets as bedding
\end{enumerate}
\par
\vspace{0.25\baselineskip}
\noindent\textbf{28.} In line 12, the phrase "under one roof" is closest in meaning to
\begin{enumerate}[label=(\Alph*),leftmargin=*,itemsep=0.2em]
\item in separate sections
\item in several shelters
\item comfortably
\item together
\end{enumerate}
\par
\vspace{0.25\baselineskip}
\noindent\textbf{29.} What was the main difference between the dwellings of Pueblo and Woodland Indians?
\begin{enumerate}[label=(\Alph*),leftmargin=*,itemsep=0.2em]
\item The Pueblos lived in permanent structures, but the Woodland Indians lived in transient shelters.
\item The Pueblos used wood in their constructions, but the Woodland Indians relied mostly on animal hides.
\item The Woodland Indians lived on flat ground, but the Pueblos lived in canyons.
\item The Woodland Indians built small shelter.
\end{enumerate}
\par
\vspace{0.25\baselineskip}
\noindent\textbf{30.} It can be inferred from the passage that Pueblo dwellings were designed to protect inhabitants from
\begin{enumerate}[label=(\Alph*),leftmargin=*,itemsep=0.2em]
\item attacks by enemies and cold winters
\item attacks by enemies and against rain water
\item wind storms and water from rain
\item wild animals, cold winters, and desert sands
\end{enumerate}
\par
\vspace{0.25\baselineskip}
\noindent\textbf{31.} The word "preclude" in line 24 is closest in meaning to
\begin{enumerate}[label=(\Alph*),leftmargin=*,itemsep=0.2em]
\item include
\item stop
\item preserve
\item conclude
\end{enumerate}
\par
\vspace{0.25\baselineskip}
\subsection*{Questions 32-42}
\begingroup
\small
\sloppy
\setlength{\parskip}{0pt}
\setlength{\parindent}{0pt}
Continents and ocean basins represent the largest identifiable bodies on Earth. On the\par
solid portions of the planet, the second most prominent features are flat plains, elevated plateaus, and large mountain ranges. In geography, the term "continent" refers to the surface of\par
continuous landmasses that together comprise about 29.2\% of the planet's surface, On the\par
other hand, another definition is prevalent in the general use of the term that deals with extensive mainlands, such as Europe or Asia, that actually represent one very large landmass.\par
Although all continents are bounded by water bodies or high mountain ranges, isolated\par
mainlands, such as Greenland and India-Pakistan areas are called subcontinents. In some\par
circles, the distinction between continents and large islands lies almost exclusively in the (10) size of a particular landmass.\par
The analysis of compression and tension in the earth's crust has determined that continental structures are composed of layers that underlie continental shelves. A great deal of\par
disagreement among geologists surrounds the issue of exactly how many layers underlie\par
each landmass because of their distinctive mineral and chemical composition. It is also quite\par
possible that the ocean floor rests on the top of unknown continents that have not yet been\par
explored. The continental crust is believed to have been formed by means of a chemical reaction when lighter materials separated from heavier ones, thus settling at various levels\par
within the crust. Assisted by the measurements of the specifics within crust formations by\par
means of monitoring earthquakes, geologists can speculate that a chemical split occurred to\par
form the atmosphere, sea water, and the crust before it solidified many centuries ago.\par
Although each continent has its special features, all consist of various combinations of\par
components that include shields, mountain belts, intracratonic basins, margins, volcanic\par
plateaus, and blockvaulted belts. The basic differences among continents lie in the proportion and the composition of these features relative to the continent size. Climatic zones have\par
a crucial effect on the weathering and formation of the surface features, soil erosion, soil\par
deposition, land formation, vegetation, and human activities.\par
Mountain belts are elongated narrow zones that have a characteristic folded sedimentary\par
organization of layers. They are typically produced during substantial crustal movements,\par
which generate faulting and mountain building. When continental margins collide, the\par
rise of a marginal edge leads to the formation of large mountain ranges, as explained by the\par
plate tectonic theory. This process also accounts for the occurrence of mountain belts in\par
ocean basins and produces evidence for the ongoing continental plate evolution.\par
\endgroup
\medskip
\noindent\textbf{32.} What does this passage mainly discuss?
\begin{enumerate}[label=(\Alph*),leftmargin=*,itemsep=0.2em]
\item Continental drift and division
\item Various definitions of the term "continent"
\item Continental structure and crust
\item Scientific analyses of continental crusts
\end{enumerate}
\par
\vspace{0.25\baselineskip}
\noindent\textbf{33.} According to the passage, how do scientists define continents?
\begin{enumerate}[label=(\Alph*),leftmargin=*,itemsep=0.2em]
\item As masses of land without divisions
\item As extensive bodies of land
\item As the largest identifiable features
\item As surficial compositions and ranges
\end{enumerate}
\par
\vspace{0.25\baselineskip}
\noindent\textbf{34.} In line 5, the word "bounded" is closest in meaning to
\begin{enumerate}[label=(\Alph*),leftmargin=*,itemsep=0.2em]
\item covered
\item convened
\item delimited
\item dominated
\end{enumerate}
\par
\vspace{0.25\baselineskip}
\noindent\textbf{35.} The author of the passage implies that the disagreement among scientists is based on the fact that
\begin{enumerate}[label=(\Alph*),leftmargin=*,itemsep=0.2em]
\item each continent has several planes and shelves
\item continents have various underlying layers of crust
\item continents undergo compression and experience tension
\item continents have different chemical makeup
\end{enumerate}
\par
\vspace{0.25\baselineskip}
\noindent\textbf{36.} The word "specifics" in line 13 is closest in meaning to
\begin{enumerate}[label=(\Alph*),leftmargin=*,itemsep=0.2em]
\item specialties
\item speculations
\item exact details
\item precise movements
\end{enumerate}
\par
\vspace{0.25\baselineskip}
\noindent\textbf{37.} The word "it" in line 15 refers to
\begin{enumerate}[label=(\Alph*),leftmargin=*,itemsep=0.2em]
\item a chemical split
\item the crust
\item the atmosphere
\item sea water
\end{enumerate}
\par
\vspace{0.25\baselineskip}
\noindent\textbf{38.} The author of the passage implies that
\begin{enumerate}[label=(\Alph*),leftmargin=*,itemsep=0.2em]
\item it is not known exactly how the continental crust was formed
\item geologists have neglected the exploration of the ocean floor
\item scientists have concentrated on monitoring earthquakes
\item the earth's atmosphere split into water and solids
\end{enumerate}
\par
\vspace{0.25\baselineskip}
\noindent\textbf{39.} According to the passage, what are the differences in the structure of continents?
\begin{enumerate}[label=(\Alph*),leftmargin=*,itemsep=0.2em]
\item The proportional size of continents to one another
\item Ratios of major components and their comparative size
\item The distinctive features of their elements
\item Climatic zones and their effect on the surface features
\end{enumerate}
\par
\vspace{0.25\baselineskip}
\noindent\textbf{40.} In line 25, the phrase "This process" refers to
\begin{enumerate}[label=(\Alph*),leftmargin=*,itemsep=0.2em]
\item continental collision
\item mountain ranges
\item the rise of margins
\item plate tectonic theory
\end{enumerate}
\par
\vspace{0.25\baselineskip}
\noindent\textbf{41.} The author of the passage implies that
\begin{enumerate}[label=(\Alph*),leftmargin=*,itemsep=0.2em]
\item the process of mountain formation has not been accounted for
\item mountain ranges on the ocean floor lead to surface mountain building
\item faulting and continental margins are parts of plate edges
\item the process of continent formation has not been completed
\end{enumerate}
\par
\vspace{0.25\baselineskip}
\noindent\textbf{42.} The word "evidence" in the last line is closest in meaning to
\begin{enumerate}[label=(\Alph*),leftmargin=*,itemsep=0.2em]
\item eventuality
\item confirmation
\item exemplification
\item challenge
\end{enumerate}
\par
\vspace{0.25\baselineskip}
\subsection*{Questions 43-50}
\begingroup
\small
\sloppy
\setlength{\parskip}{0pt}
\setlength{\parindent}{0pt}
Chicago ranks as the leading industrial and urban center in North America. Carl Sandburg called it the "City of the Big Shoulders" primarily because in the 1930s, its population\par
contained a large segment of industrial workers, the largest agricultural market, and a huge\par
airport. This poetic phrase, however, does not do justice to the city's outstanding array of\par
cultural institutions, such as the symphony orchestra and the museums of art and history.\par
The downtown busiriess district on the shore of Lake Michigan is the hub of fashionable\par
and elegant boutiques, quaint restaurants, and high-rise office buildings. Lake Shore Drive\par
extends to both the north and south ends of the city, making it one of the most spectacular\par
roadways in the state of Illinois. The Old Water Tower, dwarfed by the John Hancock Cen-\par
ter, is a must-see landmark frequented by thousands of tourists each year. Most of the Chicago lakefront is public, with spectacular beaches and wide lawns stretching along the shore\par
line.\par
Throughout its history, Europeans streamed into the city in search of jobs in steel mills\par
and factories. The large influx of population created tensions in various neighborhoods, and\par
in the 1920s, Chicago gained a reputation for violence and crime that it never lived down.\par
Nonetheless, the booming industries continued to attract new residents into the thriving\par
city, despite its notoriety.\par
The Chicago metropolitan area has undergone dramatic changes since the 1940s when\par
the suburban population almost doubled, and the number of city residents fell. Today, most\par
of the city's ethnic enclaves have faded away, but their rich heritage remains. The residents\par
take pride in impressive churches and blocks of homes constructed in the early 20th century\par
by industrious European immigrants who built the city. More than 85\% of Chicagoans\par
were born in the United States, and access to Irish, Italian, Polish, and German community\par
institutions and businesses is not as important to them as it was to their grandparents. Ital-\par
ian is no longer spoken in Little Italy, and Irish pubs have fewer Irish customers than those\par
of mixed, typically American origins.\par
\endgroup
\medskip
\noindent\textbf{43.} What does the passage mainly discuss?
\begin{enumerate}[label=(\Alph*),leftmargin=*,itemsep=0.2em]
\item Chicago's industrial and urban evolution
\item Cultural and tourist attractions in Chicago
\item The size of the city and its roadways
\item The spectacular arrays of buildings in Chicago
\end{enumerate}
\par
\vspace{0.25\baselineskip}
\noindent\textbf{44.} According to the passage, Carl Sandburg's phrase
\begin{enumerate}[label=(\Alph*),leftmargin=*,itemsep=0.2em]
\item describes the city in the best light
\item reflects the city in its entirety
\item overlooks many of Chicago's attractive features
\item gives an overview of Chicago's cultural life
\end{enumerate}
\par
\vspace{0.25\baselineskip}
\noindent\textbf{45.} The word "hub" in line 6 is closest in meaning to
\begin{enumerate}[label=(\Alph*),leftmargin=*,itemsep=0.2em]
\item hill
\item center
\item corner
\item home
\end{enumerate}
\par
\vspace{0.25\baselineskip}
\noindent\textbf{46.} The author of the passage implies that
\begin{enumerate}[label=(\Alph*),leftmargin=*,itemsep=0.2em]
\item the Old Water Tower is shorter than the John Hancock Center
\item the John Hancock Center is probably the tallest building in the city
\item the Old Water Tower and the John Hancock Center are located on the lake shore
\item the Old Water Tower and the John Hancock Center are easy to reach by one of the Illinois roadways
\end{enumerate}
\par
\vspace{0.25\baselineskip}
\noindent\textbf{47.} According to the passage, immigrants from Europe
\begin{enumerate}[label=(\Alph*),leftmargin=*,itemsep=0.2em]
\item arrived in Chicago by ships and boats
\item came to the city to enjoy its beaches
\item arrived in large numbers
\item came to Chicago to live in a large city
\end{enumerate}
\par
\vspace{0.25\baselineskip}
\noindent\textbf{48.} What changes have occurred in the city since the 1940s?
\begin{enumerate}[label=(\Alph*),leftmargin=*,itemsep=0.2em]
\item Many residents moved out to neighboring towns.
\item Its population grew rapidly.
\item Many residents forgot their ethnic heritage.
\item Its original builders moved back to Europe.
\end{enumerate}
\par
\vspace{0.25\baselineskip}
\noindent\textbf{49.} In line 20, the word "industrious" is closest in meaning to
\begin{enumerate}[label=(\Alph*),leftmargin=*,itemsep=0.2em]
\item employed in an industry
\item hard-working
\item ill-famed
\item employed in construction
\end{enumerate}
\par
\vspace{0.25\baselineskip}
\noindent\textbf{50.} According to the passage, currently most residents of Chicago
\begin{enumerate}[label=(\Alph*),leftmargin=*,itemsep=0.2em]
\item speak several languages
\item do not shop in local stores
\item do not have community institutions
\item predominantly speak English
\end{enumerate}
\par
\vspace{0.25\baselineskip}
\end{document}
